% group_meeting_math_cards.tex
% Cleanly-typeset math for the May 2026 group meeting slides.
% Compile with pdflatex; each "card" is a screenshot-friendly box.
%
% Sources:
%   docs/PNP_BV_Analytical_Simplifications.md
%   docs/clipping_conventions.md
%   docs/stern_layer_physics_and_next_steps.md
%   docs/Mangan2025_experimental_alignment.md
%   Forward/bv_solver/forms_logc.py
%   Forward/bv_solver/forms_logc_muh.py
%   Forward/bv_solver/boltzmann.py
%   writeups/WeekOfApr27/PNP Inverse Solver Revised.tex

\documentclass[11pt]{article}
\usepackage[margin=0.75in]{geometry}
\usepackage{amsmath,amssymb,bm}
\usepackage{xcolor}
\usepackage[most]{tcolorbox}
\usepackage{microtype}
\usepackage{parskip}
\usepackage{siunitx}
\usepackage{enumitem}
\setlist[itemize]{leftmargin=*,topsep=2pt,itemsep=2pt}

% Plain black-on-white card style. Title sits on a thin rule above
% the body; no colored fills or borders.
\tcbset{
  cardstyle/.style={
    enhanced,
    colback=white,
    colframe=white,
    coltitle=black,
    fonttitle=\bfseries\large,
    title=#1,
    boxrule=0pt,
    frame hidden,
    arc=0pt,
    left=4pt, right=4pt, top=4pt, bottom=4pt,
    attach boxed title to top left={xshift=0pt, yshift=0pt},
    boxed title style={
      colback=white,
      colframe=white,
      boxrule=0pt,
      arc=0pt,
      left=0pt, right=0pt, top=0pt, bottom=2pt,
    },
    after title={\par\vspace{2pt}\hrule\vspace{6pt}},
    breakable=false
  }
}

\newcommand{\VT}{V_{T}}
\newcommand{\Eeq}{E_{\mathrm{eq}}}
\newcommand{\VRHE}{V_{\mathrm{RHE}}}

\begin{document}

\begin{center}
  {\Large\bfseries Math Reference Cards --- PNP--BV Solver Updates}\\[2pt]
  {\small For the May 2026 group meeting slides. Each card is a
  self-contained block intended to be screenshotted directly into a
  slide.}
\end{center}

%% =====================================================================
\begin{tcolorbox}[cardstyle={Baseline: 4-species PNP-BV (pre-rebuild)}]
\textbf{Bulk Nernst--Planck (one equation per species $i$):}
\[
  \frac{\partial c_i}{\partial t}
   = -\nabla\!\cdot J_i,\qquad
  J_i = -D_i\!\left(\nabla c_i + \frac{z_i F}{RT}\,c_i\,\nabla\varphi\right)
\]

\textbf{Poisson:}
\[
  -\varepsilon\,\nabla^{2}\varphi \;=\; F\sum_{i} z_i\,c_i
\]

\textbf{Butler--Volmer flux BC at the electrode for reaction $j$:}
\[
  r_j \;=\; k_{0,j}\!\left[\,
  c_O^{\mathrm{surf}}\,\exp\!\Bigl(-\alpha_j\,\tfrac{n_e F\eta_j}{RT}\Bigr)
  \;-\;
  c_R^{\mathrm{surf}}\,\exp\!\Bigl((1-\alpha_j)\,\tfrac{n_e F\eta_j}{RT}\Bigr)
  \right]
\]

\textbf{Species (4-species formulation):} \(\mathrm{O_2}\),
\(\mathrm{H_2O_2}\), \(\mathrm{H^+}\), \(\mathrm{ClO_4^-}\). Charges
\(z = (0, 0, +1, -1)\). Reactions
\[
  \mathrm{R_1}:\;\mathrm{O_2} + 2\mathrm{H^+} + 2e^{-}\!\to\mathrm{H_2O_2},
  \qquad
  \mathrm{R_2}:\;\mathrm{H_2O_2} + 2\mathrm{H^+} + 2e^{-}\!\to 2\mathrm{H_2O}.
\]
\textbf{Convergence window:} \(\VRHE\in[-0.50,+0.10]\) V (cathodic
edge only).
\end{tcolorbox}

%% =====================================================================
\begin{tcolorbox}[cardstyle={Change 1: 3-species + analytic Boltzmann counterion}]
Drop \(\mathrm{ClO_4^-}\) as a dynamic Nernst--Planck unknown. Replace
its contribution to Poisson by an analytic Boltzmann profile in
local equilibrium with \(\varphi\).

\textbf{Modified Poisson source (nondimensional, ideal mode):}
\[
  -\varepsilon\,\nabla^{2}\varphi \;=\;
  F\!\!\sum_{i\in\{\mathrm{O_2,H_2O_2,H^+}\}}\!\! z_i c_i
  \;-\; F\,c_{\mathrm{ClO_4}}^{\mathrm{bulk}}\,\exp(-z_{\mathrm{ClO_4}}\,\varphi)
\]

\textbf{Why it works:} a species in zero-flux equilibrium with the
field satisfies
\(\nabla(\ln c_i + z_i\varphi) = 0\), i.e.\
\(c_i \propto \exp(-z_i\varphi)\). For an inert supporting electrolyte
this is exact in the chemically active region; the inert co-ion never
needs to be discretized.

\textbf{Convergence window after Change 1:} cold-starts in
\(\VRHE\gtrsim -0.20\) V; warm-walk fills the full
\(\VRHE\in[-0.50,+0.10]\) V grid with sub-\(0.5\%\) agreement to the
4-species reference (V24 validation, 8/8 PASS).
\end{tcolorbox}

%% =====================================================================
\begin{tcolorbox}[cardstyle={Change 2: log-concentration formulation}]
Switch primary unknowns from \(c_i\) to
\(\;\boxed{\,u_i = \ln c_i\,}\;\). Reconstruct \(c_i = \exp(u_i)\)
which is positive by construction.

\textbf{Nernst--Planck flux in log-c form:}
\[
  J_i \;=\; -D_i\,c_i\bigl(\nabla u_i + z_i\nabla\varphi\bigr)
  \;=\; -D_i\,e^{u_i}\bigl(\nabla u_i + z_i\nabla\varphi\bigr)
\]

\textbf{Backward-Euler weak form:}
\[
  \int_\Omega \frac{e^{u_i} - e^{u_i^{\mathrm{old}}}}{\Delta t}\,v\,d\Omega
  \;+\;
  \int_\Omega D_i\,e^{u_i}\bigl(\nabla u_i + z_i\nabla\varphi\bigr)\!\cdot\!\nabla v\,d\Omega
  \;=\; \text{(BV boundary terms)}
\]

\textbf{Poisson source in log-c form:}
\[
  \int_\Omega \varepsilon\,\nabla\varphi\!\cdot\!\nabla w\,d\Omega
  \;-\;\int_\Omega \rho_{\mathrm{rhs}}\!\sum_{i} z_i\,e^{u_i}\,w\,d\Omega
  \;=\; 0
\]

\textbf{Why it helps:} \(c_i = e^{u_i} > 0\) by construction; the
Debye-layer behavior \(c_i \sim c_{\mathrm{bulk}}\exp(-z_i\varphi/\VT)\)
becomes approximately linear in \(u_i\); SNES iterates can no longer
overshoot to negative concentrations.
\end{tcolorbox}

%% =====================================================================
\begin{tcolorbox}[cardstyle={Change 3: log-rate Butler--Volmer}]
Inside the BV boundary residual, build the rate \emph{additively in
log space} and exponentiate once at the end.

\textbf{Cathodic branch of reaction $j$:}
\begin{align*}
  \log r_j^{(\text{cat})}
   &= \ln k_{0,j} \;+\; u_{\mathrm{cat}(j)}
      \;+\; \sum_{f}\nu_f\bigl(u_{\mathrm{sp}(f)} - \ln c_{\mathrm{ref},f}\bigr)
      \;-\;\alpha_j\,n_e\,\eta_j/\VT \\[2pt]
  r_j^{(\text{cat})} &= \exp\!\bigl(\log r_j^{(\text{cat})}\bigr)
\end{align*}

\textbf{Anodic branch:}
\[
  \log r_j^{(\text{anod})}
   = \ln k_{0,j} + u_{\mathrm{anod}(j)}
     + (1-\alpha_j)\,n_e\,\eta_j/\VT,\qquad
  r_j^{(\text{anod})} = \exp\!\bigl(\log r_j^{(\text{anod})}\bigr)
\]

\textbf{What this fixes:} the legacy form evaluates
\(c_{\mathrm{surf},i}=\exp(\mathrm{clamp}(u_i,\pm 30))\) before
multiplying by the BV exponential. At anodic \(\VRHE\), \(u_{\mathrm{H_2O_2}}\)
becomes very negative; the clamp pins \(c_{\mathrm{surf}}\) at a
finite floor and the saturated cathodic exponential of \(\mathrm{R_2}\)
produces a phantom sink. Adding \(u_i\) inside the exponent lets
\(\log r_j \to -\infty\) smoothly, which is the physically correct
behavior.

\textbf{Convergence window after Change 3:} the production grid now
reaches \(\VRHE = +0.60\) V (was \(\le +0.20\) V).
\end{tcolorbox}

%% =====================================================================
\begin{tcolorbox}[cardstyle={Stern compact layer (Robin BC)}]
The legacy no-Stern model imposes
\(\varphi_s = \varphi_m = \varphi_{\mathrm{applied}}\) at the
electrode (Dirichlet). Real interfaces split the applied drop:
\[
  \varphi_m - \varphi_{\mathrm{bulk}}
   \;=\; \Delta\varphi_{\mathrm{Stern}} + \Delta\varphi_{\mathrm{Diffuse}}
\]

\textbf{Compact-layer capacitance law:}
\[
  \sigma \;=\; C_S\,(\varphi_m - \varphi_s)
\]

\textbf{Robin boundary contribution to the Poisson residual} (sign
matches \texttt{forms\_logc.py}):
\[
  F_{\mathrm{res}} \;\mathrel{-}=\;
   \underbrace{\widetilde{C}_S}_{\text{nondim}}\,
   (\varphi_{\mathrm{applied}}-\varphi)\,w\;d\Gamma_{\mathrm{electrode}}
\]

\textbf{Modified BV overpotential under Stern:}
\[
  \eta_j \;=\; \varphi_m - \varphi_s - \Eeq^{(j)}
\]

\textbf{Limiting behavior:}
\[
  C_S\to\infty:\;\varphi_m - \varphi_s \to 0
   \;\;(\text{recovers no-Stern Dirichlet});\qquad
  C_S\;\text{finite}:\;\varphi_s\;\text{floats; some drop in the
  compact layer}.
\]

\textbf{Why it matters:} at \(\VRHE > +0.66\) V the Boltzmann co-ion
source \(c_{\mathrm{ClO_4}}^{\mathrm{bulk}}\exp(\varphi_s)\) becomes
catastrophically large under no-Stern. Finite \(C_S\) shifts part of
the drop into the compact layer, lowering \(\varphi_s\) in the PNP
domain.

\textbf{May~4 result:} on the production stack
(3sp + Bikerman counterion + log-c-\(\mu_H\)) with the
Bikerman-consistent IC, \(C_S = 0.10\,\si{\farad\per\meter\squared}\)
gives \(15/15\) clean convergence on
\(\VRHE\in[-0.5, +1.0]\,\)V. At \(\VRHE = +0.65\,\)V the Stern layer
absorbs \(\approx 12.8\,\)V of the applied potential, keeping
\(\psi_D\) modest enough that proton supply does not underflow the
BV cathodic terms.
\end{tcolorbox}

%% =====================================================================
\begin{tcolorbox}[cardstyle={Bikerman steric correction}]
Finite ion size penalizes packing. Define the local packing fraction
\(\theta\) using ion sizes \(a_i\) (nondimensional):
\[
  \theta \;=\; 1 - \sum_{i\,\in\,\text{dyn}}\!a_i\,c_i
   \;-\;\underbrace{a_b\,c_{\mathrm{steric},b}}_{\text{counterion contribution}}
\]

\textbf{Steric chemical potential added to the species' driving force:}
\[
  \mu_i^{\mathrm{steric}} \;=\; -\ln\theta
  \quad\Longrightarrow\quad
  J_i \;=\; -D_i\,c_i\bigl(\nabla u_i + z_i\nabla\varphi
                            + \nabla\mu_i^{\mathrm{steric}}\bigr)
\]

\textbf{Pre-steric (ideal) Boltzmann counterion} --- the legacy form
already in use under Change~1 (zero-flux equilibrium in a point-ion
PNP, no finite-size penalty on the counterion itself):
\[
  c_b(x) \;=\; c_b^{\mathrm{bulk}}\,\exp\!\bigl(-z_b\,\varphi(x)\bigr)
\]
Contribution to the Poisson residual
(\texttt{boltzmann.py:add\_boltzmann\_counterion\_residual}):
\[
  F_{\mathrm{res}} \;\mathrel{-}=\;
   \rho_{\mathrm{rhs}}\,z_b\,c_b^{\mathrm{bulk}}\,
   \exp\!\bigl(-z_b\,\varphi\bigr)\,w\,d\Omega
\]
This form is \emph{unbounded} as \(\varphi\to-\infty\) (for \(z_b>0\))
or \(\varphi\to+\infty\) (for \(z_b<0\)) --- the source of the
\(\VRHE>+0.66\) V Poisson blow-up before Stern was added.

\textbf{Steric-aware analytic Boltzmann counterion} (the Bikerman
closure derived in
\texttt{steric\_analytic\_clo4\_reduction\_handoff.md}):
\[
  c_b(x) \;=\;
   c_b^{\mathrm{bulk}}\,
   \frac{\bigl(1 - \sum_{j\in\mathrm{dyn}} a_j c_j(x)\bigr)\,
         \exp(-z_b\,\varphi(x))}
        {\theta_b + a_b\,c_b^{\mathrm{bulk}}\,\exp(-z_b\,\varphi(x))}
\]
where \(\theta_b = 1 - \sum_j a_j c_j^{\mathrm{bulk}} - a_b c_b^{\mathrm{bulk}}\)
is the bulk packing fraction (must be positive). The denominator
caps \(c_b\) at \(1/a_b\) so the modified Poisson source stays
bounded everywhere.

\textbf{Limiting behavior:} \(a\to 0\) recovers the ideal Boltzmann
profile \(c \propto \exp(-z\,\varphi)\) above.
\end{tcolorbox}

%% =====================================================================
\begin{tcolorbox}[cardstyle={Debye--Boltzmann analytical IC}]
Composite outer-electroneutral profile + inner Gouy--Chapman Debye
layer. Picard-iterate the two-rate algebraic mass balance, then build
\(\varphi\) and \(u_i\) on the production mesh.

\textbf{Outer linear profiles} (ambipolar diffusion gives the proton
its effective \(2 D_H\) factor):
\[
  O_s = O_b - \tfrac{R_1}{D_O},\quad
  P_s = P_b + \tfrac{R_1 - R_2}{D_P},\quad
  H_o = H_b - \tfrac{R_1+R_2}{D_H}
\]

\textbf{Diffuse-layer Boltzmann correction:}
\[
  H_s \;=\; H_o\,\exp(-\psi_D),\qquad
  \psi_D \;=\; \varphi_{\mathrm{electrode}} - \varphi_o,\qquad
  \varphi_o \;=\; \ln\!\bigl(H_o / c_{-}^{\mathrm{bulk}}\bigr)
\]

\textbf{2$\times$2 algebraic BV system (one Picard step):}
\begin{align*}
  A_1 &= k_1 (H_s/H_{\mathrm{ref}})^2 \exp(-\alpha_1 n\eta_1), &
  B_1 &= k_1 \exp\bigl((1-\alpha_1) n\eta_1\bigr),\\
  A_2 &= k_2 (H_s/H_{\mathrm{ref}})^2 \exp(-\alpha_2 n\eta_2),
\end{align*}
\[
  \begin{bmatrix}
    1 + A_1/D_O + B_1/D_P & -B_1/D_P\\
    -A_2/D_P              & 1 + A_2/D_P
  \end{bmatrix}\!
  \begin{bmatrix} R_1 \\ R_2 \end{bmatrix}
   =
  \begin{bmatrix} A_1 O_b - B_1 P_b \\ A_2 P_b \end{bmatrix}
\]

\textbf{Gouy--Chapman composite for the potential:}
\[
  \lambda_D \;=\; \sqrt{\varepsilon/(2 H_o)},\qquad
  \psi(y) \;=\; 4\,\mathrm{atanh}\!\Bigl(\tanh(\psi_D/4)\,e^{-y/\lambda_D}\Bigr)
\]
\[
  \varphi_{\mathrm{IC}}(y) = \ln\!\bigl(H_{\mathrm{outer}}(y)/c_{-}^{\mathrm{bulk}}\bigr) + \psi(y),
  \qquad
  c_{H,\mathrm{IC}}(y) = H_{\mathrm{outer}}(y)\,e^{-\psi(y)},
  \qquad
  u_{i,\mathrm{IC}} = \ln c_{i,\mathrm{IC}}
\]

\textbf{Result:} 1.7--6.3$\times$ fewer Newton iterations at
\(\VRHE\ge +0.30\) V; unblocks \(+0.60\) V cold-start where the linear
IC fails outright. Extended further by the Bikerman-consistent IC
on the next card.
\end{tcolorbox}

%% =====================================================================
\begin{tcolorbox}[cardstyle={Bikerman-consistent IC (Option 2b, May 4)}]
Refines the Debye--Boltzmann IC for the steric (Bikerman) production
stack. Two pieces:

\textbf{(1) Composite \(\psi\) profile}
(Bazant--Kilic--Storey--Ajdari matched-asymptotic):
\[
  \nu \;=\; 2\,a_{\mathrm{Cl}}\,c_{\mathrm{Cl}}^{\mathrm{bulk}},\qquad
  \psi_{\mathrm{sat}} \;=\; \ln\!\bigl(2/\nu\bigr),\qquad
  \alpha \;=\; \sqrt{\frac{2}{\nu\,\lambda_D^{2}}\,
    \ln\!\bigl(1 + \nu(\cosh\psi_D - 1)\bigr)}
\]
\[
  y_{\mathrm{match}} \;=\; \frac{|\psi_D| - \psi_{\mathrm{sat}}}{\alpha}
\]
\[
  \psi(y) \;=\;
  \begin{cases}
    \mathrm{sgn}(\psi_D)\,(|\psi_D| - \alpha\,y)
       & 0 \le y < y_{\mathrm{match}} \\[3pt]
    \mathrm{sgn}(\psi_D)\,\psi_{\mathrm{sat}}\,
        \exp\!\bigl(-(y - y_{\mathrm{match}})/\lambda_D\bigr)
       & y \ge y_{\mathrm{match}}
  \end{cases}
\]
Falls back to the tanh/Gouy--Chapman expression of the previous card
when \(|\psi_D| \le \psi_{\mathrm{sat}}\) (linear-Debye limit).

\textbf{(2) Multispecies activity correction \(\gamma\):}
\[
  \gamma(y) \;=\;
  \frac{1}{1 + a_H\,H_{\mathrm{outer}}(y)\,(e^{-\psi(y)} - 1)
            + a_{\mathrm{Cl}}\,c_{\mathrm{Cl}}^{\mathrm{anchor}}\,
                              (e^{+\psi(y)} - 1)}
\]
(neutral species drop out automatically since \(e^{0}-1=0\)).

\textbf{Seed every species on the saturated steric manifold:}
\begin{align*}
  u_{\mathrm{O_2}}^{\mathrm{IC}}   &= \ln c_{\mathrm{O_2}}^{\mathrm{outer}}   + \ln\gamma\\
  u_{\mathrm{H_2O_2}}^{\mathrm{IC}} &= \ln c_{\mathrm{H_2O_2}}^{\mathrm{outer}} + \ln\gamma\\
  u_{H}^{\mathrm{IC}}    &= \ln H_{\mathrm{outer}} - \psi(y) + \ln\gamma\\
  \varphi^{\mathrm{IC}}  &= \ln\!\bigl(H_{\mathrm{outer}} / c_{\mathrm{Cl}}^{\mathrm{bulk}}\bigr) + \psi(y)
\end{align*}

\textbf{$\mu_H$ form:} same construction; \(\psi\) cancels exactly,
\(\mu_H^{\mathrm{IC}} = 2\ln H_{\mathrm{outer}} - \ln c_{\mathrm{Cl}}^{\mathrm{bulk}} + \ln\gamma\).

\textbf{Result (3sp+Bikerman+\(\mu_H\), \(N_y=200\), clip 50, \(u\)-clamp 100):}
\(15/15\) cold/warm convergence on
\(\VRHE\in\{-0.5,\dots,+1.0\}\,\)V with Stern \(0.10\,\)F/m\(^{2}\);
\(14/15\) without Stern. Cold ceiling \(+0.60\,\)V (was \(+0.30\,\)V),
warm-walk reach \(+1.00\,\)V (was \(+0.60\,\)V).
\end{tcolorbox}

%% =====================================================================
\begin{tcolorbox}[cardstyle={Electrochemical-potential ($\mu_H$) formulation}]
Replace the proton primary unknown \(u_H\) with the electrochemical
potential
\[
  \boxed{\,\mu_H \;=\; u_H + z_H\,\varphi\,} \qquad (z_H = +1)
\]
so that
\(\;c_H = \exp(\mu_H - z_H\varphi)\;\) is reconstructed from the muh
unknown plus the resolved field.

\textbf{Nernst--Planck flux collapses (no separate drift term):}
\[
  J_H \;=\; -D_H\,c_H\,\nabla\mu_H
\]

\textbf{Why it helps in the Debye layer:} when \(\mathrm{H^+}\) is in
local Boltzmann equilibrium, \(\nabla\mu_H \approx 0\). The primary
variable \(\mu_H\) is therefore nearly flat in \(y\) even when
\(u_H\) and \(\varphi\) each vary by tens of log-units across the
Debye layer. Newton sees a smooth unknown in exactly the regime where
\(u_H\) is most wild.

\textbf{Cancellation under the Debye-Boltzmann IC:}
\[
  \mu_H^{\mathrm{IC}}(y)
  = u_H^{\mathrm{IC}}(y) + \varphi^{\mathrm{IC}}(y)
  = \bigl[\ln H_{\mathrm{outer}} - \psi(y)\bigr]
   +\bigl[\ln(H_{\mathrm{outer}}/c_{-}^{\mathrm{bulk}}) + \psi(y)\bigr]
  = 2\ln H_{\mathrm{outer}} - \ln c_{-}^{\mathrm{bulk}}
\]
The diffuse-layer \(\psi\) cancels exactly --- a clean correctness
check.

Non-mu species (\(\mathrm{O_2}\), \(\mathrm{H_2O_2}\) --- neutral) are
unchanged from the log-c stack.
\end{tcolorbox}

%% =====================================================================
\begin{tcolorbox}[cardstyle={Local pH observable (deck-aligned)}]
The Mangan/Ruggiero deck reports IrOx-sensed local pH at the disk
surface. Cheap to derive from the existing surface \(\mathrm{H^+}\)
field (no new species, no new BC):

\[
  \mathrm{pH}_{\mathrm{local}}(t)
   \;=\; -\log_{10}\!\Bigl(c_{H,\mathrm{surf}}(t) \,/\, 1000\Bigr)
\]

with \(c_{H,\mathrm{surf}}\) in \(\si{\mol\per\meter\cubed}\) (so the
\(/1000\) converts to \(\si{\mol\per\liter}\)). In the log-c
formulation this is even cheaper:
\[
  \mathrm{pH}_{\mathrm{local}}(t)
   \;=\; -\frac{1}{\ln 10}\bigl(u_{H,\mathrm{surf}}(t) - \ln 1000\bigr).
\]

\textbf{Use:} report alongside CD/PC. Compare local-vs-bulk pH
trajectory against the Ruggiero 2022 RRDE/IrOx data.
\end{tcolorbox}

%% =====================================================================
\begin{tcolorbox}[cardstyle={RRDE transport: Levich diffusion length}]
Tie the abstract reference length \(L_{\mathrm{ref}}\) to a physical
RRDE rotation rate via the Levich diffusion-layer thickness:
\[
  L_{\mathrm{eff}} \;\approx\;
   1.61\;D_i^{1/3}\,\nu^{1/6}\,\omega^{-1/2}
\]
where \(\omega\) is the angular rotation rate
(\(\si{\radian\per\second}\)), \(\nu\) is the kinematic viscosity of
the electrolyte
(\(\si{\meter\squared\per\second}\)), and \(D_i\) is the diffusivity
of the rate-limiting species
(\(\si{\meter\squared\per\second}\)).

\textbf{Worked numbers for water at 25\,\textdegree C, O\textsubscript{2}-limited:}
\[
  D_{\mathrm{O_2}} = 1.9\times 10^{-9}\,\si{\meter\squared\per\second},\qquad
  \nu = 8.9\times 10^{-7}\,\si{\meter\squared\per\second}
\]

\(\omega\) corresponds to \(\Omega \;[\mathrm{rpm}]\) by
\(\omega = 2\pi\Omega / 60\):
\begin{center}
\begin{tabular}{rcc}
\hline
\(\Omega\,[\mathrm{rpm}]\) & \(\omega\,[\si{\radian\per\second}]\) & \(L_{\mathrm{eff}}\,[\si{\micro\meter}]\) \\
\hline
400  & 41.9  & \(\sim 86\) \\
1600 & 167.6 & \(\sim 43\) \\
2500 & 261.8 & \(\sim 35\) \\
\hline
\end{tabular}
\end{center}

\textbf{Use:} promote \(L_{\mathrm{ref}}\) from a generic
nondimensionalization scale to an \emph{experimental} parameter so
that rotation-rate sweeps in the multi-experiment FIM design
correspond to a real measurable knob.
\end{tcolorbox}

%% =====================================================================
\begin{tcolorbox}[cardstyle={Where we are vs.\ deck experiment (closure list)}]
\textbf{Aligned (physics-framework level):} PNP, Bikerman steric,
finite Stern, two-step BV (R$_1$/R$_2$), peroxide observable, three
dynamic species + analytic counterion, log-c, log-rate.

\textbf{Not yet aligned (experimental closure):}

\begin{itemize}
  \item \textbf{Electrolyte identity.} The Ruggiero study uses
        \(\mathrm{H_2SO_4 + MOH + M_2SO_4}\),
        \(M\in\{\mathrm{Li,Na,K,Cs}\}\). Our solver carries
        \(\mathrm{ClO_4^-}\) as a \emph{generic} inert surrogate.
        Need a sulfate-family balancing-anion path for deck-matched runs.
  \item \textbf{Cation specificity.} Deck story is ion-specific.
        Current solver has a generic steric size
        \(a_{\mathrm{default}} = 0.01\); needs cation-specific radii
        and an explicit (or analytic) cation species.
  \item \textbf{Local pH observable.} See card above.
  \item \textbf{RRDE transport.} See card above.
  \item \textbf{Ring channel + selectivity.} Post-processing on the
        existing peroxide-current observable.
\end{itemize}
\end{tcolorbox}

%% =====================================================================
\begin{tcolorbox}[cardstyle={Headline: voltage-window expansion}]
\begin{center}
\begin{tabular}{lccc}
\hline
& \textbf{Pre-rebuild} & \textbf{Apr 27 stack} & \textbf{May 4 stack} \\
& \textbf{(4-species)} & \textbf{(3sp+B+log-c+log-rate)} & \textbf{(+ Bikerman-IC + Stern)} \\
\hline
Cold-start window & \([-0.50,+0.10]\) & \(\gtrsim -0.20\) & \(\gtrsim -0.20\) \\
Full grid (warm-walk) & \([-0.50,+0.10]\) & \([-0.50,+0.60]\) & \([-0.50,+1.00]\) \\
Demonstrated reach
   & \(+0.10\) V
   & \(+0.60\) V
   & \(+1.00\) V (Bikerman-IC, Stern; \(15/15\)) \\
   &
   &
   & \(+2.00\) V (log-c+steric+\texttt{debye\_boltzmann}, warm-walk) \\
\(\mathrm{R_2}\) unclips at & --- & \(+0.495\) V & \(+0.495\) V \\
Local FIM \(\mathrm{cond}(F)\) & \(2.0\times 10^{11}\) & \(1.79\times 10^{7}\) & TBD \\
Weak eigenvector
   & \(\approx \log k_{0,2}\)
   & \(\approx \log k_{0,1}\)
   & TBD\\
\hline
\end{tabular}
\end{center}

\textbf{Caveat — kinetic dead zone above \(\sim +0.5\) V.}
Convergence reaches \(+2.0\) V, but converged currents drop to
machine epsilon (\(10^{-12}\)--\(10^{-16}\) mA/cm\(^{2}\)) once
\(c_{H,\mathrm{surf}} \approx \exp(-\psi_D)\) underflows the cathodic
BV terms. This is real physics, not a clip artefact: the proton
factor \((c_H/c_{\mathrm{ref}})^{2}\) (R\(_1\)) and \((c_H/c_{\mathrm{ref}})^{2}\)
(R\(_2\)) underflow at \(\psi_D > 30\). Inverse-relevant signal
remains in \(\VRHE\in[-0.5, +0.5]\) V.

\textbf{R$_2$ unclip threshold derivation.} The implementation clips
\(\eta_{\mathrm{scaled}} = (V-\Eeq)/\VT\) to \(\pm 50\) before the
\(\alpha\,n_e\) multiplication. R$_2$ becomes unclipped when
\(|V - \Eeq^{(2)}| < 50\,\VT\), i.e.\
\[
  V > \Eeq^{(2)} - 50\,\VT
   = 1.78 - 50\,(0.02569)
   = +0.495\,\text{V vs RHE}.
\]
\end{tcolorbox}

\end{document}
